\section{Exploratory Local-Model Results}
\label{sec:results}

All 128 assigned local-model generations completed without a transport failure.
The primary five-condition run contributed 80 outputs and the compact-compiler
ablation contributed 48. Raw responses, complete prompts, selected preference
records, provider token counts, finish reasons, latency, run manifests, and seeds
are retained in JSONL. The two analysis datasets also provide 64 and 32
condition-blinded response pairs, respectively, with mapping keys stored separately.

\subsection{Five-condition pilot}

The initial structured compiler (\texttt{verbose\_v1}) did not deliver the hoped-for
context advantage. Compared with history retrieval, E consumed 65.0 additional
input tokens per prompt (exploratory 95\% interval $[62.1,68.0]$), produced 52.1
additional words ($[11.6,99.1]$), and added 3.56 seconds of latency
($[1.16,6.26]$). Its reference-concept coverage difference was $-0.047$
($[-0.177,0.104]$). Compared with no personalization, E produced 34.8 fewer words
($[-75.2,-7.0]$) and returned 1.13 seconds faster ($[-2.29,-0.13]$), but required
116.0 more input tokens. The coverage difference was $0.073$
($[-0.052,0.188]$). These intervals describe this small post-hoc pilot and are not
confirmatory confidence intervals for a target population.

\subsection{Compact compiler ablation}

We replaced descriptive structured records with fixed, bounded behavioral clauses
(\texttt{compact\_v2}), reducing the example compiled profile from 447 to 158
characters. The resulting E prompts used 19.0 fewer input tokens than history
retrieval ($[-19.75,-18.25]$), thereby meeting a narrow input-context objective.
However, E produced 70.8 more words ($[28.6,117.3]$), 99.1 more output tokens
($[42.0,161.8]$), and added 3.12 seconds of latency ($[1.32,5.03]$). The rate of
length-limited termination was 0.313 higher, although its interval included zero
($[-0.063,0.750]$). Reference-concept coverage was nearly unchanged
($\Delta=0.005$, $[-0.141,0.177]$).

Against no personalization, compact E used 32.0 additional input tokens
($[31.25,32.75]$), produced 16.7 fewer words ($[-42.8,6.3]$), and returned 1.47
seconds faster ($[-2.73,-0.47]$). Its coverage difference was $-0.021$
($[-0.156,0.125]$).

\begin{table}[t]
  \centering
  \small
  \caption{Pooled paired E-minus-comparator differences across two models, four
  tasks, and two replicates. Intervals are exploratory model--task cluster
  bootstraps. Lower is better for token, word, and latency costs; reference
  coverage is a post-hoc lexical proxy.}
  \label{tab:pilot-results}
  \begin{tabular}{llrrr}
    \toprule
    Compiler & Comparator & $\Delta$ coverage & $\Delta$ input tokens &
    $\Delta$ words \\
    \midrule
    verbose\_v1 & History retrieval & $-0.047$ & $+65.0$ & $+52.1$ \\
    verbose\_v1 & No personalization & $+0.073$ & $+116.0$ & $-34.8$ \\
    compact\_v2 & History retrieval & $+0.005$ & $-19.0$ & $+70.8$ \\
    compact\_v2 & No personalization & $-0.021$ & $+32.0$ & $-16.7$ \\
    \bottomrule
  \end{tabular}
\end{table}

The strongest supported statement is therefore a systems trade-off: compact
structured selection reduced input context relative to this history baseline, but
the tested instruction semantics expanded generation enough to erase the end-to-end
efficiency benefit. Neither run establishes preference satisfaction or subjective
quality.
